%-------------------------------------------------------------------------------
%	SECTION TITLE
%-------------------------------------------------------------------------------
\cvsection{Relevante Forschungserfahrung}%
\begin{cventries}%
%---------------------------------------------------------
  \cventry%
    {PostDoc{\enskip\cdotp\enskip}Arbeitsgruppe Golling} % Advisor
    {Universit\'e de Gen\`eve -- D\'epartement de Physique Nucl\'eaire et Corpusculaire} % Institution
    {Genf, Schweiz} % Location
    {seit 2024} % Date(s)
    {%
      \begin{cvitems} % Description(s)
        \item {%
            Leitung von Studien zu \emph{Foundation Models} f\"ur die Datenanalyse in der Hochenergie-/ und Teilchenphysik%
        }%
        \item {%
            Mitarbeit an ATLAS Run~2 Legacy \ttHbb-CP-Analyse spezialisiert auf Softwareentwicklung und Sampleproduktion
        }%
        \item {%
            Runkoordinator des ATLAS Pixeldetektors im letzten Betriegsjahr des LHC in Run 3%
        }%
        \item {%
            Entwicklung von Software und Continuous Integration für das zentrale \emph{TopCPToolkit} Analyse-Framework%
        }%
      \end{cvitems}%
      ~%
    }%
%---------------------------------------------------------

  \cventry%
    {Promotion{\enskip\cdotp\enskip}Doktorvater: Prof.\ Dr.\ Arnulf Quadt} % Advisor
    {Georg-August-Universit\"at G\"ottingen -- II.\ Physikalisches Institut} % Institution
    {G\"ottingen, Deutschland} % Location
    {2020\,--\,2024} % Date(s)
    {%
      \begin{cvitems} % Description(s)
        \item {%
            Ma{\ss}gebliche Mitarbeit an ATLAS Run~2 Legacy \ttHbb-Analyse%
        }%
        \item {%
            Entwicklung und Validierung von Transformer-Modellen zur Klassifikation und Rekonstruktion von Kollisionsereignissen%
        }%
        \item {%
            C++-Softwareentwicklung für den ATLAS Pixeldetektor, die ATLAS Top Arbeitsgruppe und die Run~2 Legacy \ttHbb-Analyse%
        }%
        \item {%
            Entwicklung einer Testsuite zur automatisierten Validierung von Software und FPGA-Firmware des ATLAS Pixeldetektors%
        }%
        \item {%
            Rufbereitschaft als On-Call-Experte für den ATLAS Pixeldetektor w\"ahrend eines CERN-Aufenthalts in 2022
        }%
      \end{cvitems}%
      ~%
    }%
%---------------------------------------------------------

  % \cventry%
  %   [7cm]
  %   {Masterarbeit{\enskip\cdotp\enskip}Betreuuer: Prof.\ Didier Queloz} % Advisor
  %   {University of Cambridge -- Cavendish Astrophysics Group} % Institution
  %   {Cambridge, Vereinigtes K\"onigreich} % Location
  %   {2019} % Date(s)
  %   {%
  %     \begin{cvitems} % Description(s)
  %       \item {Masterarbeit \"uber bodengest\"utzte Entdeckung transitierender Exoplaneten mit Next-Generation Transit Survey}%
  %     \end{cvitems}%
  %     ~%
  %   }%
%-% --------------------------------------------------------

  % \cventry%
  %   {Erasmus+ Austausch{\enskip\cdotp\enskip}Betreuende: Prof.\ Dr.\ Arnulf Quadt, Prof.\ Dr.\ Baida Achkar} % Advisor
  %   {Georg-August-Universit\"at G\"ottingen -- II.\ Physikalisches Institut} % Institution
  %   {G\"ottingen, Deutschland} % Location
  %   {2018} % Date(s)
  %   {%
  %     \begin{cvitems} % Description(s)
  %       \item {Vorbereitende Studien f\"ur inklusive \ttbar Wirkungsquerschnittmessung mit ATLAS Run~2 Datensatz}%
  %     \end{cvitems}%
  %     ~%
  %   }%
%---------------------------------------------------------

  % \cventry%
  %   [7cm]
  %   {Sommerpraktikum{\enskip\cdotp\enskip}Betreuende: Dr.\ David Williams, Pratyay Poddar} % Advisor
  %   {University of Cambridge -- Hitachi Cambridge Laboratory} % Institution
  %   {Cambridge, Vereinigtes K\"onigreich} % Location
  %   {2016} % Date(s)
  %   {%
  %     \begin{cvitems} % Description(s)
  %       \item {Proof-of-Concept-Experimente zu festk\"orperbasierten Quantencomputing mit Elektronspin Qubits}%
  %     \end{cvitems}%
  %     ~%
  %   }%
%---------------------------------------------------------
\end{cventries}%
